%!TeX program = xelatex
% 使用课程提供的 SYSUReport 模板；模板作者 NorthSecond，CC BY 4.0。
\documentclass{SYSUReport}
\usepackage{array,tabularx,longtable,placeins}
\usepackage{xurl}
\usepackage{tikz}
\usetikzlibrary{arrows.meta,positioning,fit}
\setmainfont{Times New Roman}
\setCJKmainfont[AutoFakeBold=2,AutoFakeSlant=0.15]{SimSun}
\setCJKsansfont{SimHei}
\setCJKmonofont{SimSun}
\setmonofont{Consolas}
\setlength{\headheight}{16pt}
\setlength{\parskip}{0.25em}
\renewcommand{\arraystretch}{1.3}
\newcolumntype{P}[1]{>{\raggedright\arraybackslash}p{#1}}
\newcolumntype{Y}{>{\raggedright\arraybackslash}X}
\newcommand{\code}[1]{\begingroup\Urlmuskip=0mu plus 1mu\expandafter\nolinkurl\expandafter{\detokenize{#1}}\endgroup}
\newcommand{\docref}[1]{\textit{《#1》}}
\lstset{language=SQL,basicstyle=\ttfamily\small,keywordstyle=\color{myteal},
  commentstyle=\color{gray},stringstyle=\color{black},breaklines=true,
  columns=fullflexible,keepspaces=true,showstringspaces=false,
  frame=single,rulecolor=\color{gray!35},backgroundcolor=\color{gray!4},
  xleftmargin=0.4em,xrightmargin=0.4em,aboveskip=0.6em,belowskip=0.6em}
\headl{数据库原理实验}
\headc{}
\headr{第1--4周阶段报告}
\fancypagestyle{plain}{\fancyhf{}\fancyhead[L]{\fangsong 数据库原理实验}
  \fancyhead[R]{\fangsong 第1--4周阶段报告}\fancyfoot[C]{\thepage}}
\lessonTitle{数据库原理实验报告}
\reportTitle{二次元周边商品店数据库设计与实现}
\stuname{郑立峰\quad 李维政}
\stuid{24325341\quad 24325143}
\inst{计算机学院}
\major{计算机科学与技术}
\date{2026年10月4日}
\hypersetup{pdftitle={二次元周边商品店数据库设计与实现 第1至4周阶段报告},pdfauthor={郑立峰 李维政}}

\begin{document}
\hypersetup{pageanchor=false}
\cover
\thispagestyle{empty}
\clearpage
\pagenumbering{Roman}
\setcounter{page}{1}
\hypersetup{pageanchor=true}
\begingroup\small\setlength{\parskip}{0pt}\tableofcontents\endgroup
\clearpage
\pagenumbering{arabic}
\setcounter{page}{1}

\section{实验目标与阶段成果}
本阶段围绕一家小型线下二次元周边商品店，完成从经营需求到关系数据库原型的设计与实现。我们将商品、库存、会员、员工、销售、预售及采购联系起来，在 SQL Server 中建立数据库，并通过增删改查、经营统计、完整性约束和角色授权验证设计的可用性。阶段版本为 v0.1，重点是形成能够按脚本复现的数据库基础，为后续数据库应用、编程和事务实验提供统一的数据结构。

\begin{table}[htbp]
\centering\small
\begin{tabularx}{\textwidth}{P{1.3cm}P{3.1cm}Y}
\toprule 周次 & 实验主题 & 本组主要成果 \\
\midrule
第一周 & 经营需求与数据边界 & 确定单门店场景、参与者、现货与预售流程及简化规则。\\
第二周 & 关系模式设计 & 设计15张表，明确属性、域、主码、候选码、外码及业务样例。\\
第三周 & DDL与数据修改 & 完成建库建表、样例装载及商品、库存、销售订单 CRUD。\\
第四周 & 查询、视图、约束与授权 & 完成9条经营查询、4个视图、完整性正反例和3种员工角色权限测试。\\
\bottomrule
\end{tabularx}
\caption{第1至4周实验任务与阶段成果}
\end{table}

实验在 Windows 环境中使用 Microsoft SQL Server 和 SQL Server Management Studio（SSMS）完成，截图所示本地实例版本为17.0.1135，数据库名为\code{AnimeGoodsStoreDB}。我们在 SSMS 中按编号执行00--07脚本，并由组员检查运行结果、业务含义和文档说明。报告围绕总体结构与典型案例展开，各表完整字段定义和完整 SQL 保留在协同仓库中，便于进一步查阅和复现。

\section{经营场景与业务需求}
\subsection{场景选择与数据边界}
我们选择徽章、亚克力立牌、挂件、色纸及玩偶等商品构成的线下周边店。相比线上商城或连锁经营，单门店能够清楚体现销售、库存和员工职责，又不必引入配送、跨店调拨等额外流程。商品的作品/IP、具体角色和预售活动为常规零售场景提供了特色，也为后续销量分析与补货决策留下扩展空间。

进入数据库的信息包括商品类别与资料、IP及角色、当前库存、会员、员工、供应商、销售订单与明细、预售活动与预订、采购订单与明细。它们分别服务于商品识别、交易记录、需求汇总和采购执行。普通非会员顾客不建立个人档案，其订单的会员字段为空；供应商虽然需要存档，但不登录店铺系统。线上商城、物流、售后退换货、多门店、复杂营销和自动推荐暂不纳入本阶段，避免业务范围超出课程进度。

\subsection{角色职责与关键业务约定}
\begin{table}[htbp]
\centering\small
\begin{tabularx}{\textwidth}{P{2.4cm}Y}
\toprule 角色或参与者 & 职责与定位 \\
\midrule
店长/管理员 & 管理商品、员工与供应商；开展预售，确定采购及额外备货；查看经营统计。\\
销售/收银店员 & 查询商品和可售库存，处理会员与销售订单，确认预售提货。\\
库存管理员 & 到货验收、库存盘点、库存维护及登记采购执行结果。\\
会员 & 购买现货、参加预售、取消预订及提货；后续通过应用访问本人数据。\\
普通顾客、供应商 & 分别参与现货购买和供货，不直接拥有数据库操作账号。\\
\bottomrule
\end{tabularx}
\caption{业务参与者与职责划分}
\end{table}

小组讨论后确定三项简化约定：一个商品最多属于一个 IP、对应一个具体角色，群像商品不指定具体角色；一次预订整体取消或整体提货，不支持部分处理；一条预订最多对应一笔预售提货订单，不合并多个活动的预订结账。这些约定保留了预售的业务特点，同时使关系结构和实验操作保持可解释、可实施。

\subsection{现货、预售与补货流程}
现货业务从查询可售库存开始，随后创建订单和明细、确认数量与成交价、完成结账，并维护库存。库存达到或低于警戒值时，由库存管理员发现补货需求，店长确定采购量，供应商到货后再验收入库。

预售则先由店长设置活动时间、预售价和订金，会员预订并支付订金，截止后统一汇总需求和确定采购。截止前取消退还订金，截止后仍可自行取消但订金不退；到货后有效预订对应商品被预留，会员在期限内补尾款并整体提货。采购后取消、逾期未提货或额外备货产生的剩余商品进入普通现货销售。因此，\textbf{预售与现货是同一商品的两种销售方式，不是两个固定商品类别}。

\begin{figure}[htbp]
\centering
\begin{tikzpicture}[font=\small,>=Stealth,
 box/.style={draw=myteal,rounded corners,fill=myteal!4,align=center,text width=3.5cm,minimum height=0.95cm},
 link/.style={->,thick,color=myteal}]
\node[box] (a) at (0,0) {开展预售\quad 会员预订\\记录订金与预订数量};
\node[box] (b) at (4.3,0) {截止后汇总有效需求\\店长确定额外备货};
\node[box] (c) at (8.6,0) {统一采购\quad 到货验收\\形成实际库存与预留};
\node[box] (d) at (8.6,-1.8) {会员补尾款并整体提货\\记录销售订单与明细};
\node[box] (e) at (4.3,-1.8) {剩余商品作为现货\\取消、逾期或额外备货};
\node[box] (f) at (0,-1.8) {现货销售与库存检查\\按需补货并验收入库};
\draw[link] (a)--(b);\draw[link] (b)--(c);\draw[link] (c)--(d);
\draw[link] (c.south west)--(e.north east);\draw[link] (e)--(f);
\end{tikzpicture}
\caption{预售与现货之间的经营流程关系}
\end{figure}

这些流程用于指导表结构和样例组织。本阶段已实现数据存储、查询和声明式约束；自动库存联动、跨表退款判断及完整状态迁移计划结合后续存储过程和事务实现。详细业务说明见\docref{第一周 业务设计.md}。

\FloatBarrier
\section{关系模式与典型设计}
\subsection{实体及数据表总体结构}
15张表按业务职责组织，而不是对每一种操作单独建表。商品基础资料提供分类维度，销售与采购分别采用“订单头＋明细”结构，预售活动与预订连接商品和会员，库存保存当前经营状态。

\begin{table}[htbp]
\centering\small
\begin{tabularx}{\textwidth}{P{2.2cm}YP{3.3cm}}
\toprule 模块 & 数据表 & 核心用途 \\
\midrule
商品基础资料 & \code{ProductCategory}、\code{AnimeIP}、\code{CharacterInfo}、\code{Product} & 分类、作品、角色与商品资料。\\
人员与供应商 & \code{Member}、\code{BusinessRole}、\code{Employee}、\code{Supplier} & 会员、经营职责、员工和供货来源。\\
库存 & \code{Inventory} & 实际库存、预留库存与补货警戒值。\\
销售 & \code{SalesOrder}、\code{SalesOrderItem} & 现货与预售提货的统一交易记录。\\
预售 & \code{PresaleActivity}、\code{PresaleReservation} & 活动安排、需求、订金与预订状态。\\
采购 & \code{PurchaseOrder}、\code{PurchaseOrderItem} & 普通补货与预售采购及验收数量。\\
\bottomrule
\end{tabularx}
\caption{实体及数据表清单}
\end{table}

\begin{figure}[htbp]
\centering
\begin{tikzpicture}[font=\footnotesize,>=Stealth,
 box/.style={draw=myteal,rounded corners,align=center,fill=myteal!4,minimum width=3.4cm,minimum height=0.65cm},
 link/.style={->,color=myteal,thick}]
\node[box] (cat) at (0,0) {ProductCategory / AnimeIP};
\node[box] (ch) at (5,0) {CharacterInfo};
\node[box] (p) at (0,-1.15) {Product};
\node[box] (i) at (5,-1.15) {Inventory};
\draw[link] (p)--(cat);\draw[link] (ch)--(cat);\draw[link] (p)--(ch);\draw[link] (i)--(p);
\node[box] (pa) at (0,-2.5) {PresaleActivity};
\node[box] (pr) at (5,-2.5) {PresaleReservation};
\node[box] (m) at (10,-2.5) {Member};
\draw[link] (pa)--(p);\draw[link] (pr)--(pa);\draw[link] (pr)--(m);
\node[box] (soi) at (0,-3.85) {SalesOrderItem};
\node[box] (so) at (5,-3.85) {SalesOrder};
\node[box] (emp) at (10,-3.85) {Employee / BusinessRole};
\draw[link] (soi)--(so);\draw[link] (so)--(pr);\draw[link] (so)--(m);\draw[link] (so)--(emp);
\draw[link] (soi.west)--++(-0.4,0)|-(p.west);
\node[box] (poi) at (0,-5.2) {PurchaseOrderItem};
\node[box] (po) at (5,-5.2) {PurchaseOrder};
\node[box] (sup) at (10,-5.2) {Supplier};
\draw[link] (poi)--(po);\draw[link] (po)--(sup);\draw[link] (po)--(emp);
\draw[link] (poi.west)--++(-0.7,0)|-(p.west);
\end{tikzpicture}
\caption{主要关系概览，箭头由引用方指向被引用方；基础资料和员工资料合并展示}
\label{fig:relations}
\end{figure}

图\ref{fig:relations}突出主要数据关联：销售和采购通过明细连接商品；预订连接活动与会员；销售订单还可关联预订。员工资料内部由\code{Employee.role_id}引用\code{BusinessRole}，预售创建者、采购经办与验收员工也关联员工表。各表完整字段、关系模式和样例元组详见\docref{第二周 数据库设计.md}，以下选择典型结构解释设计理由。

\subsection{商品设计与作品角色一致性}
\code{Product}使用内部编号\code{product_id}作为主码，业务编号\code{product_code}设置唯一约束作为候选码。二者分别承担表间引用与业务识别，商品名称不要求唯一，避免不同款式同名时无法录入。

\begin{table}[htbp]
\centering\small
\begin{tabularx}{\textwidth}{P{3.2cm}P{2.8cm}Y}
\toprule 典型字段 & 类型及约束 & 设计含义 \\
\midrule
\code{product_id} & INT，PK，正整数 & 稳定的内部标识。\\
\code{product_code} & VARCHAR(20)，UNIQUE & 商品业务编号。\\
\code{category_id} & INT，非空外码 & 每件商品必须有类别。\\
\code{ip_id}、\code{character_id} & INT，可空外码 & 区分无明确IP、群像及具体角色商品。\\
\code{sale_price} & DECIMAL(10,2)，大于0 & 当前售价，以定点数保存金额。\\
\code{product_status} & VARCHAR(20)，CHECK & 限定 PENDING、ON\_SALE、STOPPED。\\
\bottomrule
\end{tabularx}
\caption{商品表的典型字段设计}
\end{table}

中文名称使用 NVARCHAR，业务编号使用 VARCHAR，金额使用 DECIMAL，业务时间使用 DATETIME2；空值用于表达确实不存在或尚未确定的信息，不用随意的字符串占位。原神群像徽章可以只填 IP，而甘雨立牌同时关联 IP 与角色。角色不为空时必须指定 IP，再由复合外键保证两者属于同一条角色记录：
\par\addvspace{0.6\baselineskip}\noindent\begin{minipage}{\linewidth}
\begin{lstlisting}
CHECK (character_id IS NULL OR ip_id IS NOT NULL),
FOREIGN KEY (character_id, ip_id)
    REFERENCES dbo.CharacterInfo(character_id, ip_id)
\end{lstlisting}
\end{minipage}\par\addvspace{0.3\baselineskip}
仅设置两个独立外键只能保证“IP存在”和“角色存在”，不能排除“原神＋三月七”的错误组合。为支持复合引用，角色表另设\code{UNIQUE(character_id, ip_id)}；该组合包含已有主码，是辅助引用的唯一约束，不是新的最小候选码。角色名称真正的业务候选码是\code{(ip_id, character_name)}。

\subsection{销售订单与明细的统一结构}
\code{SalesOrder}记录订单编号、销售类型、会员、收银员工、预订关联及完成时间；\code{SalesOrderItem}记录商品、数量和实际成交单价，其联合主码为\code{(order_id, product_id)}。订单头与明细分离，使一笔现货订单能够容纳多种商品，同时避免每条明细重复保存会员和收银员工。

普通顾客的\code{member_id}可以为空。现货订单为 SPOT，预订关联为空；预售提货订单为 PRESALE，必须关联会员和预订。我们用\code{(reservation_id, member_id)}复合外键保证提货会员与预订会员一致，并用过滤唯一索引限制一条非空预订最多形成一笔订单：
\par\addvspace{0.6\baselineskip}\noindent\begin{minipage}{\linewidth}
\begin{lstlisting}
CREATE UNIQUE INDEX UX_SalesOrder_ReservationID
ON dbo.SalesOrder(reservation_id)
WHERE reservation_id IS NOT NULL;
\end{lstlisting}
\end{minipage}\par\addvspace{0.3\baselineskip}
这样既限制重复提货订单，也允许多笔不关联预订的普通现货订单。\code{order_no}是订单业务候选码；上述索引表达的是非空预订关联的条件唯一性。

成交单价必须保存在明细中，不能直接取商品当前售价，否则后续调价会改变历史销售统计。订单总金额通过\code{SUM(quantity * unit_price)}计算，不再重复保存。样例 SO004 包含甘雨立牌1件和派蒙徽章2件，总金额为\(68+2\times28=124\)元。预订表也不重复保存最终成交价和提货时间，它们分别由销售明细和订单完成时间表达。

\subsection{库存状态与预售决策数据}
\code{Inventory.product_id}同时是主码和商品外码，保证同一商品最多存在一条当前库存记录；商品与库存的结构关系为\(1:0..1\)。正常经营时为商品初始化库存，表结构本身不强制每个商品都有库存记录。

库存只保存实际数量、已到货商品的预留数量和补货警戒值，并约束：
\[
0\leq\text{reserved\_qty}\leq\text{on\_hand\_qty},\qquad
\text{available\_qty}=\text{on\_hand\_qty}-\text{reserved\_qty}.
\]
例如在库5件、预留2件时，可售库存为3件；若其中一件预留因逾期释放，在库量不变，可售库存变为4件。这种结构无需另建“预售库存”和“现货库存”两套记录。

预售决策字段采用不同原则：\code{decision_demand_qty}保存采购决策发生时的需求快照，\code{extra_stock_qty}保存店长的额外备货选择，\code{planned_purchase_qty}保存最终采购计划。历史快照不会随取消或提货而改写，因此与查询实时计算的 ACTIVE 需求可以不同。可计算的当前状态尽量不重复保存，但具有历史决策意义的信息需要保留。

\FloatBarrier
\section{数据库建立与数据修改实验}
\subsection{执行顺序与样例数据}
完整执行链为：
\begin{center}
\code{00}建库/复用\(\rightarrow\)\code{01}重建表\(\rightarrow\)\code{02}样例装载\(\rightarrow\)\code{03} CRUD\\[0.4em]
\(\rightarrow\)\code{04}查询\(\rightarrow\)\code{05}视图\(\rightarrow\)\code{06}约束验证\(\rightarrow\)\code{07}权限验证
\end{center}
建表按照先被引用、后引用的顺序进行，重建时先删除视图，再按相反依赖顺序删除业务表。00脚本在数据库不存在时创建，存在时复用；01脚本负责清空并重建本项目结构。这里的重复运行是课程样例的确定性重建，不是保留业务数据的增量更新。05使用\code{CREATE OR ALTER VIEW}，06和07对补充约束、角色及测试用户设置存在性判断。

样例数据覆盖全部15张表，其中商品6条、库存6条、会员5条、销售订单6条、销售明细8条、预售活动2条、预订4条、采购订单3条、采购明细5条。它们体现了非会员购买、会员多商品订单、预售提货、采购后取消和开放中预售等情况。我们没有只插入彼此独立的元组，而是让商品、会员、员工、预订、销售和采购记录形成可供查询的关联。

\begin{figure}[htbp]
\centering
\includegraphics[width=\textwidth]{setup.png}
\caption{SSMS中的数据库创建成功消息}
\end{figure}

\subsection{典型 CRUD 流程与结果}
03脚本使用9001系列独立编号。商品实验先新增测试商品与库存，查询确认后修改售价和商品名称；库存由10件改为15件，警戒值由3改为5，并同步修改\code{updated_at}。查询使用相同 WHERE 条件确认目标，避免无条件更新。图\ref{fig:crud}展示了这些操作的前后变化。

\begin{figure}[htbp]
\centering\includegraphics[width=\textwidth]{crud.png}
\caption{商品信息修改与库存由10件改为15件的运行结果}
\label{fig:crud}
\end{figure}

订单实验先创建 CREATED 订单及明细，将数量从1改为2，按历史单价68元得到明细金额136元；确认明细后再将订单更新为 COMPLETED 并记录完成时间。该顺序对应“先确认交易内容，再完成结账”，避免演示中在订单完成后继续调整明细。

删除时先删除订单明细再删除订单，商品删除前先清理对应库存，体现外键对依赖关系的保护。实验结束后清理临时编号。经人工运行和检查，商品、库存与订单相关增删改查均符合预期，正式样例可继续用于后续查询。完整步骤见\docref{第三周 脚本说明.md}及\code{03_crud.sql}。

\FloatBarrier
\section{经营查询与统计视图}
\subsection{查询覆盖与统计口径}
查询围绕经营问题组织，涵盖连接、聚合、分组、HAVING和子查询。销量与销售额只统计 COMPLETED 订单，金额按历史成交单价计算；订单详情则保留订单类型和状态，便于查看业务记录。

\begin{table}[htbp]
\centering\small
\begin{tabularx}{\textwidth}{P{1.1cm}YP{4.1cm}}
\toprule 编号 & 业务问题 & 主要 SQL 方法 \\
\midrule
Q1 & 谁购买了什么，由谁收银，明细金额是多少？ & 多表 INNER JOIN、会员 LEFT JOIN。\\
Q2 & 各商品的累计销量与销售额是多少？ & LEFT JOIN、CASE、SUM、GROUP BY。\\
Q3、Q4 & 各IP、各商品类别的销售表现如何？ & 分类维度连接与分组聚合。\\
Q5 & 每位会员完成几笔订单、累计消费多少？ & LEFT JOIN、COUNT DISTINCT、SUM。\\
Q6 & 哪些商品需要关注补货？ & 可售库存计算与阈值筛选。\\
Q7 & 各预售活动当前还有多少有效需求？ & 按 ACTIVE 预订聚合，与快照对照。\\
Q8 & 哪些商品累计销量不少于2件？ & HAVING 筛选聚合结果。\\
Q9 & 哪些消费会员的累计消费高于平均值？ & 派生表、子查询与 AVG。\\
\bottomrule
\end{tabularx}
\caption{9条经营查询的总体覆盖}
\end{table}

\subsection{连接与聚合的典型分析}
Q1连接订单、明细、商品和员工，对会员使用 LEFT JOIN。原因是非会员购买时\code{member_id=NULL}，若使用 INNER JOIN，SO001、SO006 等正常订单会消失。多商品订单出现多行明细则是正常结果，不能将明细行数直接当成订单数。

Q5从会员出发，将完成状态条件放入 LEFT JOIN 的 ON 子句，使未消费会员也保留。核心实现如下：
\par\addvspace{0.6\baselineskip}\noindent\begin{minipage}{\linewidth}
\begin{lstlisting}
SELECT m.member_no,
       COUNT(DISTINCT so.order_id) AS completed_order_count,
       SUM(COALESCE(soi.quantity * soi.unit_price, 0))
           AS total_consumption
FROM dbo.Member AS m
LEFT JOIN dbo.SalesOrder AS so
  ON m.member_id = so.member_id
 AND so.order_status = 'COMPLETED'
LEFT JOIN dbo.SalesOrderItem AS soi
  ON so.order_id = soi.order_id
GROUP BY m.member_no;
\end{lstlisting}
\end{minipage}\par\addvspace{0.3\baselineskip}
林晓的 SO003、SO004 共2笔订单，但连接后有3条明细；COUNT DISTINCT 得到正确订单数，累计消费为\(129+68+2\times28=253\)元。陈雨消费90元，周琪消费74元，另外两位会员为0。Q9的平均对象是\textbf{有实际消费的会员}，平均累计消费为\((253+90+74)/3=139\)元，因此筛选出林晓；这与对全部注册会员求平均是不同的业务口径。

商品、IP和类别统计提供了另一组一致性检查：样例累计售出10件，销售额530元。原神相关商品7件、327元，星穹铁道相关商品2件、164元，初音未来相关商品1件、39元，分类总额与商品总额一致。运行结果见图\ref{fig:query}。

\begin{figure}[htbp]
\centering\includegraphics[width=0.92\textwidth]{query.png}
\caption{商品与IP维度的销量、销售额统计结果}
\label{fig:query}
\end{figure}

Q7区分当前仍需履约的需求和当时采购决策：PS001的历史需求为2件，额外备货1件、计划采购3件；一条预订已提货、另一条已取消，当前 ACTIVE 需求为0。PS002两条有效预订分别为1件、2件，当前需求为3件，采购快照尚未填写。保留两类信息，才能解释经营过程中的需求变化。Q8则对分组后的\code{SUM(quantity)}使用 HAVING，得到 P001、P002、P004 三种销量不少于2件的商品。经人工核验，各查询结果与样例业务含义一致。

\subsection{视图封装与复用}
我们将4类高频查询封装为视图：\code{vw_OrderDetail}用于订单详情，\code{vw_ProductSalesSummary}用于商品销量及销售额，\code{vw_InventoryStatus}用于库存状态，\code{vw_MemberConsumptionSummary}用于会员消费统计。视图复用基础表数据和计算规则，也为不同员工提供适合职责的查询入口。

库存状态视图实时计算可售库存，先判断数量为0的 OUT\_OF\_STOCK，再判断达到或低于警戒值的 LOW\_STOCK，其余为 NORMAL。样例 P003实际库存2件、无预留、警戒值3件，被标记为 LOW\_STOCK，其他商品为 NORMAL，见图\ref{fig:stock}。

\begin{figure}[htbp]
\centering\includegraphics[width=\textwidth]{stock.png}
\caption{库存状态视图对可售库存与补货警戒值的判断}
\label{fig:stock}
\end{figure}

视图查询结果经人工检查与原查询一致。详细查询与视图设计见\docref{第四周 设计说明.md}，完整实现分别位于\code{04_query.sql}和\code{05_view.sql}。

\FloatBarrier
\section{完整性约束与正反例验证}
\subsection{约束总体设计与测试方法}
建表阶段使用 PRIMARY KEY 保证内部标识唯一，UNIQUE 保护业务编号等候选码，FOREIGN KEY 维护引用关系，NOT NULL 和 DEFAULT 表达必填及初始状态，CHECK 限制金额、数量、枚举状态及同一行内字段关系。第四周补充“已退订金不超过实付订金”，并用2组合法正例和14组非法反例检查实际效果。

\begin{table}[htbp]
\centering\small
\begin{tabularx}{\textwidth}{P{1.8cm}YP{3.2cm}}
\toprule 测试 & 验证内容 & 预期机制 \\
\midrule
正例P1、P2 & 合法商品与库存；合法订金预订。 & 成功插入，随后回滚。\\
T1--T4 & 负售价、预留超过在库、重复商品编号、不存在的商品类别。 & CHECK、UNIQUE、外键。\\
T5、T6 & IP与角色组合错误；退款超过订金。 & 复合外键、CHECK。\\
T7--T9 & 订金高于预售价、提货期限过早、取消缺少取消时间。 & 金额、时间、状态字段约束。\\
T10、T11 & 提货会员与预订不一致；完成订单缺少交易时间。 & 复合外键、状态时间约束。\\
T12--T14 & 采购状态与验收信息不一致；预计或实际到货早于下单。 & 验收信息与时间约束。\\
\bottomrule
\end{tabularx}
\caption{完整性正反例测试的覆盖范围}
\end{table}

测试使用独立编号，在事务中执行并回滚，避免污染正式样例。反例通过 TRY/CATCH 捕获错误，同时核对错误号和预期约束名：CHECK与外键冲突通常检查547，重复商品编号检查2601或2627。只有预期规则拒绝输入才输出 PASS，不能把其他 SQL 错误误认为约束成功。这里的事务用于隔离测试数据，不代表已经实现销售和采购的完整业务事务。

\subsection{典型正反例及结果分析}
\textbf{合法商品与库存。} 正例P1插入合法商品，并设置实际库存10件、预留3件、警戒值5件。商品分类、IP及角色引用均有效，库存满足数量关系，数据库允许插入。图\ref{fig:constraint}展示正例以及部分反例的运行消息。

\textbf{IP与角色不一致。} T5使用\code{ip_id=1}（原神）和\code{character_id=3}（三月七）。两个编号分别存在，但组合不属于同一角色记录，因此\code{FK_Product_Character_IP}拒绝插入。该例说明参照完整性不仅是“引用存在”，还可以表达关联身份的一致性。

\textbf{退款超过实付订金。} T6设置\code{deposit_paid=20}、\code{deposit_refunded=30}，由\code{CK_PresaleReservation_RefundNotExceedPaid}拒绝。这是同一行金额关系，可直接用 CHECK 强制；截止前后是否退款则需读取活动截止时间，安排在后续业务过程中处理。

\textbf{预售订单会员不匹配。} T10引用属于会员2的预订R002，却把订单会员填写为1。\code{FK_SalesOrder_ReservationMember}对两字段组合检查并拒绝插入，避免为其他会员的预订建立提货订单。

\begin{figure}[htbp]
\centering\includegraphics[width=\textwidth]{constraint.png}
\caption{合法数据通过与非法商品价格被约束拒绝的运行消息}
\label{fig:constraint}
\end{figure}

小组已对完整测试进行人工审验，合法数据成功，T1--T14均因相应预期规则被拒绝，结果符合预期；回滚后未留下临时测试记录。完整用例与约束名称见\code{06_constraint.sql}。这些约束能保护已覆盖的数据关系，但状态取值合法与状态迁移合法并不相同，后者仍需后续流程控制。

\FloatBarrier
\section{角色权限与正常越权测试}
\subsection{最小权限与角色映射}
\code{BusinessRole}保存员工的现实经营职责，SQL Server数据库角色真正控制访问，两者不能等同。我们创建\code{dbrole_sales_clerk}、\code{dbrole_inventory_manager}和\code{dbrole_store_manager}三个角色，按对象和字段逐项授权，不直接赋予 db\_owner 等全权限角色。

\begin{table}[htbp]
\centering\small
\begin{tabularx}{\textwidth}{P{2.0cm}YYY}
\toprule 数据或操作 & 销售店员 & 库存管理员 & 店长 \\
\midrule
商品资料 & 查询。 & 查询。 & 增删改查。\\
会员资料 & 查询、登记、修改。 & 不授权。 & 查询。\\
销售订单 & 新增；修改状态、时间；按订单ID定位。 & 不授权。 & 查询。\\
销售明细 & 新增；修改数量和单价。 & 不授权。 & 查询。\\
库存 & 通过状态视图查询。 & 查询；修改数量、警戒值和更新时间。 & 查询；修改警戒值。\\
预售 & 查询活动、预订；修改预订状态。 & 不授权。 & 管理活动；查询预订。\\
采购 & 不授权。 & 查询；修改验收信息与到货数量。 & 管理采购订单及明细。\\
员工与供应商 & 不授权。 & 不授权。 & 管理基础资料。\\
\bottomrule
\end{tabularx}
\caption{主要对象与操作的角色权限概览}
\end{table}

视图授权使员工不必获得所有底层表权限即可使用常见结果，例如销售店员通过库存状态视图查询可售数量，通过订单详情视图查看销售记录。列级 UPDATE 则限定其可以修改的字段。会员的“只能访问本人记录”属于后续应用身份和行范围控制，本阶段不直接为会员授予基础表权限；供应商也不建立登录角色。

\subsection{典型权限验证}
为避免额外建立真实服务器登录，我们创建3个 WITHOUT LOGIN数据库用户并加入对应角色，以\code{EXECUTE AS USER}切换上下文，测试后 REVERT，修改类测试回滚。

正常测试P1由销售店员创建订单及明细，然后将 CREATED 更新为 COMPLETED；P2由库存管理员更新补货警戒值和时间；P3由店长修改商品售价。越权测试N1尝试让销售店员修改采购订单，N2让库存管理员读取会员数据，N3让销售店员直接修改实际库存。测试同时覆盖“该做的操作能完成”和“不属于职责的操作被拒绝”。

\par\addvspace{0.6\baselineskip}\noindent\begin{minipage}{\linewidth}
\begin{lstlisting}
EXECUTE AS USER = 'u_test_inventory';
EXEC sys.sp_executesql N'SELECT * FROM dbo.Member;';
-- Permission test code catches error 229 and checks Member.
REVERT;
\end{lstlisting}
\end{minipage}\par\addvspace{0.3\baselineskip}
N2对应库存管理员无权访问会员个人资料，数据库返回229权限拒绝错误。N1、N3执行越权修改同样会被访问控制阻止；UPDATE中的目标定位也需要读取相关列，实际检查过程中可能首先遇到 SELECT 权限拒绝。测试脚本据错误号与目标对象确认操作确实被权限系统阻止。经人工审验，3组正常测试与3组越权测试均符合预期，见图\ref{fig:permissions}。

\begin{figure}[htbp]
\centering\includegraphics[width=\textwidth]{permissions.png}
\caption{三种角色的正常操作成功，越权操作被数据库拒绝}
\label{fig:permissions}
\end{figure}

权限实验中有一项有代表性的修正：起初销售店员已有 INSERT和状态、时间列的 UPDATE 权限，但P1在\code{WHERE order_id=9301}定位订单时仍缺少 SELECT 权限。我们没有扩大为整个订单表的读取权限，而是补充：
\par\addvspace{0.6\baselineskip}\noindent\begin{minipage}{\linewidth}
\begin{lstlisting}
GRANT SELECT (order_id)
ON OBJECT::dbo.SalesOrder TO dbrole_sales_clerk;
\end{lstlisting}
\end{minipage}\par\addvspace{0.3\baselineskip}
修正后P1正常完成。该问题说明最小权限不是单纯减少授权，还要考虑完整 SQL 操作所需的读取条件。完整授权和测试见\code{07_role.sql}。

\FloatBarrier
\section{小组协作与 AI 辅助下的独立判断}
\subsection{成员分工与协同方式}
小组成员为郑立峰（24325341）、李维政（24325143），均来自计算机学院计算机科学与技术专业。我们先共同讨论场景和实验安排，再采用实现、复核、修订的协作方式，避免仅由代码编写者判断自己的结果。

\begin{table}[htbp]
\centering\small
\begin{tabularx}{\textwidth}{P{1.8cm}YY}
\toprule 环节 & 郑立峰 & 李维政 \\
\midrule
需求与安排 & 共同讨论场景、业务边界和阶段任务。 & 共同讨论场景、业务边界和阶段任务。\\
实现与检查 & 协同AI实现和修订第1至4周业务产出。 & 检查实现方向，对已有产出进行完整二次人工审核。\\
阶段整理 & 根据审核结果多轮修订代码与设计文档。 & 撰写README与协作文档。\\
报告与收尾 & 负责阶段报告编写与整理。 & 参与报告审阅、核对设计与实验表述。\\
\bottomrule
\end{tabularx}
\caption{第1至4周小组分工安排}
\end{table}

小组讨论记录保留了场景规模、预售取消、提货期限和采购规则等选择过程。各周文档明确设计，SQL落实结构与操作，README组织复现入口，AI对话记录保留建议及人工修改，分工文档说明责任。这些材料与阶段报告各有用途，共同组成可理解、可运行的阶段提交。

\subsection{AI建议、人工选择与验证}
AI主要辅助提出候选方案、整理数据结构、生成候选 SQL和审查不一致。我们对业务含义、复杂度及实现方式作出选择，而不是将每条建议直接纳入项目。

\textbf{预售的建模修正。} 早期方案曾把商品分为预售和现货两类。人工讨论认为同一商品可能先预售、到货后再现货销售，因此改为两种销售业务，并独立建立预售活动表。取消是否需要店铺批准、截止后是否退订金、逾期如何处理，也由小组确定后再落实到业务设计。

\textbf{复杂度与数据冗余的取舍。} 对多IP、多角色的中间表建议，我们选择暂不纳入，采用可空的商品IP与角色字段；对订单总金额、重复提货时间等字段，则采纳删除冗余的建议。这里既不是追求表越多越好，也不是一律少存字段，而是区分当前可计算值与历史交易、决策信息。

\textbf{实现与运行检查。} AI协助生成建表顺序、查询和测试框架后，我们在本地执行脚本，对结果作人工核对。出现销售正常操作失败时，进一步分析 SQL 所需读权限，并选择只补\code{order_id}列的 SELECT。修改有明确问题依据，而不是为让测试通过而直接授予全表权限。

\subsection{整合审查与改进}
阶段收尾时，我们将审查发现的问题分成“当前应修正”和“后续实现”两类。当前修正包括：调整CRUD顺序，先确认明细再完成订单；将反例判断由“执行失败”改为核对预期错误原因；补齐状态、时间及预售身份一致性约束；统一空库与已有库上的完整执行链。跨表库存联动、退款与状态机则写入待实现清单，避免文档将业务目标提前表述成已经自动实现。

会员消费聚合曾出现“消除了 NULL 值”的警告。我们将\code{COALESCE(SUM(...),0)}调整为\code{SUM(COALESCE(...,0))}，先将无明细的金额转换为0再聚合，保持业务含义并使输出更清晰。最终代码经组员复核和运行检验，阶段文档同步修订。上述过程体现了AI辅助提案、人工判断、实际运行及迭代修正之间的分工，详细记录见\docref{小组讨论议题.md}及仓库的“第1-4周ai对话”目录。

\section{实验总结与后续工作}
本阶段将经营需求转化为15张关联数据表，在 SQL Server 中完成建库、样例装载和 CRUD，再通过9条查询、4个视图、完整性正反例和角色权限验证形成数据库原型。商品、订单、预售与采购不再是孤立信息，能够支持订单追踪、销售统计、库存检查和需求分析。

实验中最重要的认识是：设计应从业务语义出发。现货与预售的区分影响活动和订单结构；历史成交价与采购快照的保留服务于后续解释；复合外键保护的不只是编号存在，还包括身份匹配。与此同时，SQL运行成功并不意味着业务解释正确，连接方式、聚合口径和权限条件都需要结合样例和职责逐项判断。

协作方面，方向讨论、二次人工审核和多轮修订帮助我们发现单人实现容易忽略的问题。AI提高了候选设计和 SQL 整理效率，但对建议的取舍、业务边界的确定及运行结果的解释仍由小组承担。我们通过具体问题的修改理解了设计，而不是只获得一组能够执行的脚本。

下一阶段将继续使用现有经营场景，结合ER建模与规范化反思关系结构，并推进数据库应用。后续数据库编程与事务阶段重点实现销售完成时的库存检查和扣减、采购验收入库、预售预留与释放、订金退款判断以及合法状态迁移。会员本人数据的访问由应用身份和数据范围控制落实。相关规则已集中记录在\docref{后续阶段待实现业务规则.md}，为后续任务提供清楚的衔接。

\subsection*{相关材料索引}
\addcontentsline{toc}{subsection}{相关材料索引}
课程要求见\docref{第一周任务讲解.docx}至\docref{第四周任务讲解.docx}。协同仓库\code{Database_exp}中，\code{README.md}提供项目范围和复现方法，“第1-4周业务产出”目录保存各周设计文档及00--07 SQL脚本，\code{results}保存运行截图，\docref{小组分工.md}和\docref{小组讨论议题.md}记录协作与AI使用过程。本报告图中的运行结果摘自既有SSMS截图，便于阅读而选取相应结果区域。

\end{document}
